\scanpage{12}
\margin{-}\margin{-}CHAPITRE I

ARITHM\handwritten{\'}{E}TIQUE DES CORPS VALU\handwritten{\'}{E}S .

\S~1 . \deleted{\unsure{C}}Valeurs absolues et valuations .

Rappelons qu'étant donné un corps $K$ , on appelle \uline{valeur absolue} sur $K$ une application $f$ de $K$ dans $R_{+}$ telle que $f(xy)=f(x)f(y)$ , $f(x+y)\leqslant f(x)+f(y)$ ,et que $f(x)=0$ entraîne $x=0$ ; la distance $f(x-y)$ définit alors sur $K$ une topologie compatible avec sa structure de corps .

Soit maintenant $v$ une valuation de $K$ dont le groupe des ordres est archimédien , c'est-à-dire isomorphe à un sous-groupe additif de $R$ . La fonction numérique $e^{-v(x)}$ est une valeur absolue sur $K$ ; on dit qu'une telle valeur absolue est \uline{valuative} .

PROPOSITION 1 .- \uline{Pour qu'une valeur absolue $f$ d'un corps} $K$ \uline{soit valuative} , \uline{il faut et il suffit que l'une ou l'autre des conditions suivantes soit vérifiée} :

a) $f(x+y)\leqslant \operatorname{Max}(f(x),f(y))$ ;

b) $f(n.1)\leqslant 1$ \uline{pour tout entier naturel $n$} .

C'est évident pour a) . a) entraîne b) car $f(1)=1$ . Réciproquement la formule du binôme montre que $f((a+b)^n)\leqslant (n+1)\operatorname{Max}(f(a)^n,f(b)^n)$ , d'où a) en faisant tendre $n$ vers l'infini .

COROLLAIRE \deleted{1} .- \uline{Toute valeur absolue d'un corps de caractéristique} $p\neq0$ \uline{est valuative} .

En effet , tout élément $n.1$ est une racine de l'unité .

On dit (\uline{Top.gén}.,chap.IX,\S~3\deleted{\unsure{1|4}},n\textsuperscript{o}2) que deux valeurs absolues $f,g$ de $K$ sont \uline{équivalentes} s'il existe un nombre réel $s$ tel que $f(x)=g(x)^s$ .

PROPOSITION 2 .- \uline{Toute valeur absolue $f$ de} $Q$ \uline{est équivalente à une valeur absolue} $p$-\uline{adique , ou à la valeur absolue ordinaire} .

Ceci a déjà été vu dans le cas où $f$ est valuative (chap. des valuations) .

Sinon , soit $n$ un entier tel que $f(n)>1$ . Soient $a$ et $b$ des entiers \inserted{\handwritten{$>1$}} ; écrivons $\deleted{\unsure{a^{b}=a^{0}|a^{b}=a_{0}}}\;a^{h}=a_{0}+a_{1}b+\ldots+a_{q}b^{q}$ , où $0\leqslant a_{i}\leqslant b-1$ . On en déduit

% archive sheet number: 11
